\documentclass[10pt]{article}
\usepackage{icmlko}

\icmlmeta{IP 파운데이션 월드모델}{254}{아홉 중 하나만 그렇다}{제 학습이 해로운 도메인은 도서뿐이다 --- 그리고 그 사실로는 판을 못 바꾼다}

\begin{document}

\twocolumn[
\icmlheader

\begin{abstract}
노트 253은 도서의 제 학습 자료가 도서를 해친다는 것을 쟀다(80건 $\to$ 40건에서
$\rho$ 0.321 $\to$ 0.360). 남은 물음은 \textbf{그게 도서만의 일인가}였다.
여덟 도메인에 같은 검사를 했다 --- 제 학습을 무작위로 절반 덜어 내고 유보
$\rho$ 를 본다. \textbf{도서만 오른다}($+$0.042). 나머지 일곱은 평평하거나
내려가고, 애니는 크게 내려간다($-$0.038, $t{=}-4.02$). 학습 크기로는 설명이
안 된다 --- spearman(학습 크기, 줄인 이득) $=$ $-$0.02 다. 도서는
\emph{작아서} 이상한 게 아니라 \textbf{다른 모집단이라} 이상하다(노트 253).
그런데 이 노트의 더 중요한 절반은 \textbf{그래서 무엇을 할 수 있나}이고,
답은 \textbf{아무것도}다. 도서를 고쳐도 판은 $+$0.0026 움직인다(몫 6.2\%
$\times$ 0.042) --- 노트 245가 그은 0.010 선 한참 아래라 안쪽으로도 유보로도
정할 수 없다. \textbf{옳다고 아는 변경을 채택할 수 없는 자리}이고, 이 실험실이
그 자리를 만난 것은 처음이 아니다(노트 216 · 224 · 234 · 241 · 247). 다만
이번엔 방향이 반대다 --- 앞의 다섯은 ``판이 오르는데 고를 수 없어서'' 거부했고,
이번은 ``고를 수 있는데 판이 안 움직여서'' 못 채택한다. 제품 층에서는 답이
다르다: 도구가 파는 것은 판 $\rho$ 가 아니라 \emph{한 도메인의 예측}이므로
(노트 247), \textbf{도서를 쓰는 사용자에게는 도서 학습 행을 빼는 것이 맞다.}
판정 단위와 배포 단위가 다르다는 것을 처음으로 정면으로 적는다.
\end{abstract}
]

\section{여덟 곳에 같은 검사를 했다}

노트 253의 검사는 한 줄이다 --- 제 학습을 덜어 내고 유보에서 재라. 도서에서
$-$0.019/ln 이 나왔다. 같은 것을 학습 30건 이상인 여덟 도메인에 다 했다
(무작위 절반, 씨앗 둘).

\begin{figure}[t]
\centering
\includegraphics[width=\columnwidth]{figs/half.pdf}
\caption{제 학습을 절반으로 줄였을 때. 오른쪽으로 가면 제 자료가 해롭다는 뜻.}
\end{figure}

\begin{center}\small
\begin{tabular}{lrrrr}
\hline
도메인 & 학습 & 지금 & 절반만 & 차 \\
\hline
\textbf{도서} & 80 & 0.321 & 0.363 & $\mathbf{+0.042}$ \\
게임 & 259 & 0.640 & 0.639 & $-$0.001 \\
모바일 & 1,559 & 0.537 & 0.535 & $-$0.002 \\
세계애니 & 2,648 & 0.520 & 0.511 & $-$0.009 \\
펀딩 & 320 & 0.236 & 0.220 & $-$0.016 \\
웹툰 & 2,106 & 0.381 & 0.349 & $-$0.032 \\
애니 & 1,467 & 0.487 & 0.449 & $\mathbf{-0.038}$ \\
아이돌 & 54 & 0.154 & 0.109 & $-$0.045 \\
\hline
\end{tabular}
\end{center}

\textbf{도서 하나만 오른다.} 애니는 $t{=}-4.02$ 로 뚜렷하게 내려간다 ---
제 자료가 확실히 쓸모 있다는 뜻이다.

크기 이야기가 아니다. spearman(학습 크기, 줄인 이득) $=$ $-$0.02 로 아무
관계가 없다. 학습 54건인 아이돌은 $-$0.045 이고 2,648건인 세계애니는
$-$0.009 다. 도서는 \emph{작아서}가 아니라 \textbf{다른 모집단이라}
이상하다 --- 노트 253이 짚은 대로 ``2020년 이전 베스트셀러''는 ``지금 잘
팔리는 옛 책''이다.

\paragraph{한 가지 흠} 표의 짝 SE 와 $t$ 는 \emph{씨앗 하나}로 잰 것이고
``차''는 씨앗 둘의 평균이라, 둘이 어긋나는 줄이 있다(웹툰은 차가
$-$0.032 인데 씨앗 0 만 보면 $+$이다). 도서와 애니는 0 에서 충분히 멀어
읽을 수 있지만 가운데 다섯 줄의 $t$ 는 안 읽는 게 맞다.

\section{그래서 무엇을 할 수 있나 --- 아무것도}

도서가 이상하다는 것은 알았다. 고치면 얼마인가.

\[ 0.042 \times 0.062 = +0.0026. \]

노트 245가 그은 0.010 선의 4분의 1 이다. 그 자리는 안쪽 신뢰도가 0.34 이고
(노트 244) 유보로 정하면 엿보기다(노트 245). \textbf{옳다고 아는 변경인데
채택할 수 없다.}

이 실험실이 이 벽을 만난 것은 여섯 번째다. 그런데 앞의 다섯과 방향이 반대다.

\begin{center}\small
\begin{tabular}{lll}
\hline
노트 & 무엇 & 왜 못 하나 \\
\hline
216 & 과제별 혼합 가중 & 판 $+$0.0050 인데 고를 수 없다 \\
224 & 웹툰 회차 간격 & 판 $+$0.0294 인데 사후다 \\
234 & 공유 시리즈 채움 & 값이 틀렸다 \\
241 & 도메인별 기울기 & 판 $+$0.0102 인데 안쪽이 거꾸로 \\
247 & 도메인 안 백분위 & 상쇄돼 판이 안 움직인다 \\
\textbf{254} & \textbf{도서 학습 빼기} & \textbf{옳은데 판이 안 움직인다} \\
\hline
\end{tabular}
\end{center}

앞의 넷은 ``판이 오르는데 근거가 없어서'' 거부했다. 254는
\textbf{근거가 있는데 판이 안 움직여서} 못 채택한다. 판 $\rho$ 가 순위
기준인 한 이 둘은 똑같이 ``안 함''으로 끝난다.

\section{판정 단위와 배포 단위가 다르다}

그런데 ``안 함''이 정말 옳은가. 노트 247이 지나가듯 적은 문장을 정면으로
꺼낼 때가 됐다.

\begin{quote}\small
도구가 파는 것은 판 $\rho$ 가 아니라 \textbf{한 도메인의 예측}이다.
\end{quote}

도서를 쓰는 사용자에게 도서 $\rho$ 는 0.321 이 아니라 0.363 일 수 있다.
그 사용자에게 ``판 전체로는 0.0026 이라 안 했습니다''는 답이 아니다.
반대로 애니 사용자에게 같은 처방을 하면 $-$0.038 이다.

\textbf{판정은 판 $\rho$ 로 하되 배포는 도메인별로 한다} --- 이것이 지금
자료가 지지하는 형태다. 근거는 노트 252의 표에 이미 있었다: 다섯 도메인은
여지가 음수라 손대면 안 되고, 도서는 제 자료가 해롭다. \textbf{하나의
전역 설정이 아홉 도메인에 다 맞을 이유가 없다.}

다만 이것을 실제로 하려면 도메인별 결정을 \emph{도메인별로} 정당화해야
하고, 도서 하나의 $+$0.042 는 짝 SE 0.024 로 $t{=}1.49$ 다. \textbf{아직
근거가 모자란다.} 이 노트는 규약을 바꾸자는 제안이 아니라, 바꿔야 할
자리를 처음으로 정확히 짚는 것이다.

\paragraph{남은 것}
\begin{center}\small
\begin{tabular}{ll}
\hline
갈래 & 왜 \\
\hline
도메인별 배포 & 판정과 배포를 가르면 무엇이 필요한가 \\
애니 $-$0.038 & 제 자료가 제일 값진 도메인. 무엇이 다른가 \\
도서 $t{=}1.49$ & 근거가 모자란다. 씨앗을 늘리면 되나 \\
웹툰 & 씨앗에 부호가 갈린다. 절반 줄이기가 불안정 \\
\hline
\end{tabular}
\end{center}

\paragraph{규약} \textbf{``판이 안 움직인다''와 ``효과가 없다''를 가른다.}
몫이 6\%인 도메인은 제 $\rho$ 가 0.04 움직여도 판에 0.0026 이다. 노트 247의
상쇄와 다른 이유로 같은 일이 생긴다 --- 그때는 부호가 엇갈려서, 이번엔
몫이 작아서.

\end{document}
